\documentclass[letterpaper,journal]{IEEEtran}

\usepackage{amsmath}
\usepackage{booktabs}
\usepackage{algorithm}
\usepackage{algpseudocode}
\usepackage{url}

\renewcommand{\thetable}{S\arabic{table}}
\renewcommand{\theequation}{S\arabic{equation}}
\renewcommand{\thealgorithm}{S\arabic{algorithm}}

\begin{document}

\title{Supplementary Material: Baseline Reproduction and Implementation Scope}
\author{Patrick Enenche and Hyowon Kim}
\maketitle

\section{Purpose and Scope}

This supplement documents the reproduction checks and implementation scope of
the K1-RSA, K3-RSA, and FSA-RSA(SSFM) baselines used in the accompanying
manuscript. The published values are taken from Tables~IV and~V of
Srivastava and Singh~\cite{Srivastava2025SSFM}; the reproduced values are
ten-seed means from the principal simulation data.

The reference study expresses traffic as an average per-node load, whereas
our simulator generates one network-wide request stream. Its Table~IV was
therefore used as the low-load reproduction check when fixing the
topology-specific load mappings. These mappings were then held constant for
every scheme and nominal load index.

\section{Low-Load Reproduction Check}

Table~\ref{tab:supp_low_load} reports the published and reproduced request
blocking probabilities at the lowest operating point. The reproduced
FSA-RSA(SSFM) values preserve the near-zero blocking regime of the reference
study.

\begin{table*}[t]
\centering
\caption{Published/reproduced low-load request blocking probability.}
\label{tab:supp_low_load}
\scriptsize
\setlength{\tabcolsep}{3.2pt}
\begin{tabular}{lccc}
\toprule
\textbf{Scheme} & \textbf{German} & \textbf{USNET} & \textbf{TI} \\
\midrule
K1-RSA
& \(4.20{\times}10^{-4}/4.80{\times}10^{-4}\)
& \(7.70{\times}10^{-4}/7.34{\times}10^{-4}\)
& \(5.00{\times}10^{-2}/4.38{\times}10^{-2}\) \\
K3-RSA
& \(1.50{\times}10^{-4}/1.99{\times}10^{-4}\)
& \(1.90{\times}10^{-4}/3.32{\times}10^{-4}\)
& \(6.00{\times}10^{-3}/2.60{\times}10^{-2}\) \\
FSA-RSA(SSFM)
& \(0/0\)
& \(0/0\)
& \(6.50{\times}10^{-4}/3.46{\times}10^{-4}\) \\
\bottomrule
\end{tabular}
\end{table*}

\section{Multi-Load Source-Style Comparison}

For the sole purpose of reproducing the presentation of Table~V in the
reference study, let \(\beta\in\{\mathrm{K1},\mathrm{K3}\}\) denote a
baseline, \(T\) a topology, and \(\mathcal{Q}\) a metric. Let
\(\mathcal{Q}_{\mathrm{SSFM},T}(\rho)\) and
\(\mathcal{Q}_{\beta,T}(\rho)\) denote the corresponding seed-averaged
values at nominal load index \(\rho\). Because the reference study does
not state the Table~V aggregation formula explicitly, we use the following
source-style signed change and report the calculation transparently:
\begin{equation}
\delta_{\mathcal{Q}}^{\mathrm{src}}(\beta,T)
=
20
\sum_{\rho\in\{5,10,15,20,25\}}
\frac{
\mathcal{Q}_{\mathrm{SSFM},T}(\rho)
-
\mathcal{Q}_{\beta,T}(\rho)
}{
\mathcal{Q}_{\beta,T}(\rho)
}.
\label{eq:supp_source_style_change}
\end{equation}
Negative request-blocking values indicate lower blocking under
FSA-RSA(SSFM), whereas positive frequency-slot-utilization and throughput
values indicate increases. Equation~\eqref{eq:supp_source_style_change} is
used only for comparison with the reference table. The principal ACA-RSA
analysis instead computes relative changes between load-averaged metric
values, thereby avoiding undefined load-specific ratios when blocking is
zero.

\begin{table*}[t]
\centering
\caption{Published/reproduced signed percentage changes for FSA-RSA(SSFM)
relative to K1-RSA and K3-RSA. The source-study LFM corresponds to the FDM
reported in our simulator.}
\label{tab:supp_table5}
\scriptsize
\setlength{\tabcolsep}{5pt}
\begin{tabular}{llrrrr}
\toprule
\textbf{Topology} & \textbf{Baseline} &
\textbf{RBP} & \textbf{FSU} & \textbf{LFM/FDM} & \textbf{Throughput} \\
\midrule
German & K1-RSA & \(-48.30/-47.23\) & \(3.00/2.83\)
& \(5.10/4.78\) & \(4.30/4.25\) \\
German & K3-RSA & \(-31.70/-29.93\) & \(-3.90/-3.91\)
& \(-1.20/-1.26\) & \(0.30/0.04\) \\
USNET & K1-RSA & \(-52.70/-53.08\) & \(6.10/6.35\)
& \(8.20/7.90\) & \(5.40/5.49\) \\
USNET & K3-RSA & \(-33.10/-36.56\) & \(-2.00/-0.24\)
& \(0.47/2.19\) & \(0.40/0.55\) \\
TI & K1-RSA & \(-50.00/-51.33\) & \(14.00/13.60\)
& \(11.70/10.56\) & \(17.20/17.14\) \\
TI & K3-RSA & \(-25.00/-39.69\) & \(-5.00/2.97\)
& \(-4.70/2.16\) & \(2.20/7.18\) \\
\bottomrule
\end{tabular}
\end{table*}

For K1-RSA, all twelve topology--metric differences are within
\(1.33\) percentage points of the published values. The K3-RSA differences
are larger for TI. The reference study describes K3-RSA as a fixed-route
scheme but does not specify how the three routes are generated or ordered.
Our implementation is therefore documented below and treated as a
common-simulator implementation rather than an exact reproduction.

\section{Baseline Procedures}

\begin{algorithm}[t]
\caption{Implemented K1-RSA}
\label{alg:supp_k1}
\begin{algorithmic}[1]
\Require Request \(R_i=(s_i,d_i,g_i)\); precomputed shortest
physical-distance path \(p_{s_i,d_i}^{(1)}\)
\State Apply First-Fit to find \(g_i\) contiguous common-free slots on
\(p_{s_i,d_i}^{(1)}\)
\If{a block is available}
    \State Allocate the block and accept \(R_i\)
\Else
    \State Block \(R_i\)
\EndIf
\end{algorithmic}
\end{algorithm}

\begin{algorithm}[t]
\caption{Implemented K3-RSA}
\label{alg:supp_k3}
\begin{algorithmic}[1]
\Require Request \(R_i=(s_i,d_i,g_i)\); three precomputed shortest
physical-distance paths \(p_{s_i,d_i}^{(1)},p_{s_i,d_i}^{(2)},
p_{s_i,d_i}^{(3)}\)
\For{\(k=1,2,3\) in increasing path-cost order}
    \State Apply First-Fit to find \(g_i\) contiguous common-free slots on
    \(p_{s_i,d_i}^{(k)}\)
    \If{a block is available}
        \State Allocate the block, accept \(R_i\), and stop
    \EndIf
\EndFor
\State Block \(R_i\) if all three attempts fail
\end{algorithmic}
\end{algorithm}

\begin{algorithm}[t]
\caption{Implemented FSA-RSA(SSFM)}
\label{alg:supp_ssfm}
\begin{algorithmic}[1]
\Require Request \(R_i=(s_i,d_i,g_i)\); current link-slot state
\For{each directed link \(e\)}
    \State Compute the current SSFM value and set
    \(\omega_e=1+\mathrm{SSFM}_e\)
\EndFor
\State Compute one minimum-weight path \(p_i\) using Dijkstra's algorithm
\State Apply First-Fit to find \(g_i\) contiguous common-free slots on
\(p_i\)
\If{a block is available}
    \State Allocate the block and accept \(R_i\)
\Else
    \State Block \(R_i\)
\EndIf
\end{algorithmic}
\end{algorithm}

All three schemes use the same event engine, request realizations, spectrum
state, and First-Fit implementation. K1-RSA and K3-RSA use route tables fixed
before event processing, whereas FSA-RSA(SSFM) recomputes its link weights and
route at each arrival.

\section{Reproducibility Files}

The accompanying package includes the processed v31 data subset and a
standard-library Python script that regenerates Tables~\ref{tab:supp_low_load}
and~\ref{tab:supp_table5}. The source-study PDF is not redistributed.

\begin{thebibliography}{1}
\bibitem{Srivastava2025SSFM}
R. Srivastava and Y. N. Singh, ``A novel fragmentation metric and
fragmentation-aware adaptive routing and spectrum allocation algorithm in
elastic optical network,'' \emph{Optical Fiber Technology}, vol.~94,
Art.~no.~104318, 2025.
\end{thebibliography}

\end{document}

